\documentclass{ximera}

\begin{document}

    \title{Lecture 4}
    
    \begin{abstract}  
    Outline: \\
    - cardinality \\
    - convergence of sequences
\end{abstract}

\maketitle

\begin{definition}{(1.4.6)}
    Let $A,B$ be subsets of some space.  A function $f:A\to B$ is one-to-one (1-1) if $a_1 \neq a_2$ in $A$ implies $f(a_1)\neq f(a_2)$ in $B$.  The function is onto if given any $b \in B$, there exists an $a \in A$ such that $f(a) = b$.
\end{definition}

We say two sets have the same cardinality if there is a one-to-one and onto (also called bijective) function mapping one to the other.

\begin{itemize}
\item A set is called \textit{countable} if it has the same cardinality as $\mathbb{N}$.
\item The rationals $\mathbb{Q}$ are countable (the proof is in your text, but a more intuitive version of the proof is the Cantor Snake, video posted on our Canvas page)
\item The real numbers are uncountable, i.e. there is no bijective mapping of $\mathbb{N}$ to $\mathbb{R}$; the proof can be done multiple ways, one of which is by contradiction using the nested interval property.  One of the most famous proofs is Cantor's diagonalization argument.
\item Your text has lovely proofs showing that any interval of the real numbers (really a characteristic one, but the proof can be generalized) has the same cardinality as the real numbers.  They are very straightforward, I recommend reading them.
\end{itemize}

This means there are at least two different sizes of infinity, one I can count and one I cannot count.  There are actually more, and the end of chapter 1 has a really nice exposition on what these different sizes of infinity are as well as the math history behind this investigation into cardinality.

\newpage

Chapter 2: Sequences \& Series

\begin{definition}{(2.2.1)}
    A sequence is a function whose domain is $\mathbb{N}$.
\end{definition}

We could write this as $f : \mathbb{N} \to A$ for some set $A$.  We often shortcut this with the notation $\{a_n \; : n \in \mathbb{N}\}$, or $\{a_n\}_{n=1}^\infty$, or even shorter, $\{a_n\}$.  In this course, $A \subseteq \mathbb{R}$.

\begin{exercise}
    Come up with sequences that, as $n \to \infty$,
    \begin{enumerate}
        \item Goes to $0$
        \item Goes to $\infty$
        \item Does not converge to anything (e.g., alternates)
    \end{enumerate}
\end{exercise}

Let's formalize this notion of ``goes to."

\begin{definition}
    A sequence $\{a_n\}$ converges to a real number $a$ if, for all $\epsilon > 0$, there exists $N \in \mathbb{N}$ such that $|a_n - a| < \epsilon$ for all $n \geq N$.
\end{definition}

\begin{exercise}
Let us consider the sequence $\left\{\frac{1}n\right\}$. What do you conjecture this sequence converges to?  Draw a number line and see if you can apply the definition for $\epsilon =\frac{1}{\sqrt2}$ and $\epsilon = \frac{2}{101}$.
\end{exercise}

% Depending on how strong the groups are, have them consequently prove the following theorem; this is what we did in class today, every group save one was able to prove this.

\begin{theorem}
    The sequence $\left\{\frac1n\right\}$ converges to $0$.
\end{theorem}
\begin{proof}
    Let $\epsilon > 0$ be a real number.  By the Archimedean Property(b), for every $\epsilon > 0$ there exists an $N$ such that $\frac{1}N < \epsilon$.  Since $\frac{1}{n} \leq \frac{1}N$ for any $n \geq N$, the proof is complete.
\end{proof}

Method of proof for sequence convergence:
\begin{itemize}
    \item Choose $\epsilon > 0$ arbitrary (for direct proofs, which nearly all convergence proofs are, this \textit{must} remain arbitrary)
    \item Find a choice of $N \in \mathbb{N}$ (most of the work is here) such that $|a_N - a| < \epsilon$
    \item Assume $n\geq N$
    \item Show that $|a_n -a| < \epsilon$, which should follow straightforwardly from the steps you have done before.
\end{itemize}

This is the usual definition we often work with, but as we have been talking about there is importance of thinking about properties in different ways different ways...

\begin{definition}{(2.2.4)}
    Given $a \in \mathbb{R}$ and $\epsilon > 0$, the set \[V_\epsilon(a) := \{ x \in \mathbb{R}: |x-a| < \epsilon \} \]
    is called an $\epsilon$-neighborhood of $a$.
\end{definition}

%\textcolor{red}{Have them use this definition to recast Definition 3 with you}

\begin{definition}
    A sequence $\{a_n\}$ converges to $a$ if for all $\epsilon > 0$ and corresponding $\epsilon$-neighborhood $V_\epsilon(a)$, there exists a point in the sequence after which all the elements are in $V_\epsilon(a)$.
\end{definition}


\end{document}